\documentclass[10pt]{article}
\usepackage{compresr}
\cprSetHeader{TypeScript · LangChain Integration}

\begin{document}

\cprTitleBlock{Compresr SDK · TypeScript}{LangChain Integration}{%
    Auto-compress tool outputs in your LangChain.js agents. Cut tokens
    30--70\%, keep answer quality, ship in one line of middleware.}

% ===========================================================================
\section{Why it matters}

LangChain.js agents burn most of their tokens on \emph{tool outputs} —
search results, scraped pages, database rows — that the model rereads on
every turn. Compresr replaces each raw output with a query-aware,
compressed version \emph{before} it enters the message state, so your KV
cache, your prompt, and your bill all shrink. Drop it in, keep your agent
exactly as is.

% ===========================================================================
\section{Install}

\begin{shcode}
npm install @compresr/sdk
export COMPRESR_API_KEY=cmp_...
\end{shcode}

\noindent Node.js 20+. Dual ESM + CJS build; the
\cprInlineCode{@compresr/sdk/integrations/langchain} subpath needs
\cprInlineCode{langchain>=1.0.0} as a peer (and any
\cprInlineCode{@langchain/*} packages you already use).

% ===========================================================================
\section{Drop-in agent middleware}

One import, one middleware entry. The agent code below is unchanged from
the LangChain.js docs except for the highlighted line.

\begin{tscode}
import { createAgent } from 'langchain';
import { compresrToolMiddleware } from '@compresr/sdk/integrations/langchain';

const agent = createAgent({
  model: 'openai:gpt-4o-mini',
  tools: [webSearch],
  middleware: [
    compresrToolMiddleware({
      apiKey: process.env.COMPRESR_API_KEY!,
      queryArg: 'query',            // use args.query as the intent
      targetCompressionRatio: 0.5,  // remove ~50% of tokens
      minTokens: 200,               // skip short outputs
    }),
  ],
});
\end{tscode}

\noindent That's it. The agent still sees every tool result, just shorter
and focused on the current query. Fail-open by default: if the Compresr
API is unreachable, the original output passes through and your agent
keeps running.

% ===========================================================================
\section{Wrap a single tool}

Prefer to compress just one noisy tool (e.g. a web scraper) inside an LCEL
chain or LangGraph node? Use the wrapper — name, description, and schema
are preserved.

\begin{tscode}
import { wrapToolWithCompression } from '@compresr/sdk/integrations/langchain';

const compressedSearch = wrapToolWithCompression(rawSearch, {
  apiKey: process.env.COMPRESR_API_KEY!,
  queryArg: 'query',
  targetCompressionRatio: 0.5,
});
\end{tscode}

\noindent Also available: \cprInlineCode{compresrSummarizationMiddleware}
(condenses conversation history at a token threshold) and
\cprInlineCode{compresrPromptMiddleware} (enforces a hard prompt budget).
For retriever-side compression, drop \cprInlineCode{CompresrExtractor}
into a \cprInlineCode{ContextualCompressionRetriever}:

\begin{tscode}
import { CompresrExtractor } from '@compresr/sdk/integrations/langchain';
import { ContextualCompressionRetriever } from 'langchain/retrievers/contextual_compression';

const retriever = new ContextualCompressionRetriever({
  baseRetriever: vectorStore.asRetriever({ k: 10 }),
  baseCompressor: new CompresrExtractor({
    apiKey: process.env.COMPRESR_API_KEY!,
    targetCompressionRatio: 0.5,
  }),
});
\end{tscode}

% ===========================================================================
\section{Key parameters}

\renewcommand{\arraystretch}{1.15}
\begin{longtable}{@{}p{4.6cm} p{9.4cm}@{}}
\toprule
\textbf{Parameter} & \textbf{Default --- Description} \\
\midrule
\endhead

\cprInlineCode{apiKey} & \emph{required} --- your Compresr API key. \\

\cprInlineCode{targetCompressionRatio} & \cprInlineCode{0.5} --- values in
$[0,1]$ remove that fraction of tokens; values $>1$ are an Nx factor
(e.g.\ \cprInlineCode{3} for $3\times$ shorter). \\

\cprInlineCode{minTokens} & \cprInlineCode{200} --- skip outputs shorter
than this. Keeps small results untouched and fast. \\

\cprInlineCode{queryArg} & \emph{undefined} --- name of the tool argument
to use as the compression query (e.g.\ \cprInlineCode{'query'}). \\

\cprInlineCode{query} & \emph{undefined} --- static query string, used
when every call shares the same intent. \\

\cprInlineCode{queryExtractor} & \emph{undefined} --- custom function to
derive a query. Middleware receives
\cprInlineCode{(toolCall, messages) => string | undefined};
\cprInlineCode{wrapToolWithCompression} receives
\cprInlineCode{(args) => string | undefined}. Return
\cprInlineCode{undefined} to fall back to the default. \\

\cprInlineCode{compressionModel} & \cprInlineCode{'latte\_v1'} ---
query-aware default. Switch to \cprInlineCode{'latte\_v2'} for the
faster variant. \\

\cprInlineCode{coarse} & \emph{undefined} --- \cprInlineCode{true} for
paragraph-level (faster), \cprInlineCode{false} for token-level
(tighter). \\

\cprInlineCode{allowTools} / \cprInlineCode{ignoreTools} & \emph{undefined}
--- limit compression to (or exclude) specific tool names. \\

\cprInlineCode{onError} & \cprInlineCode{'passthrough'} --- fail-open
behaviour. Set to \cprInlineCode{'raise'} to fail-closed. \\

\bottomrule
\end{longtable}

\vspace{0.4em}
\noindent\textbf{Get started:} \cprInlineCode{npm install @compresr/sdk}
\quad\textbullet\quad Docs and pricing at \texttt{compresr.ai}.

\end{document}
